\uuid{UmYA}
\titre{Partie entière}
\auteur{bodin ; hanani ; mzari}
\organisation{exo7}
\licence{CC BY-NC-SA 4.0}
\datecreate{2020-01-20}
\niveau{L1}
\module{Analyse}
\chapitre{Propriétés de R}
\sousChapitre{Propriétés des nombres réels}
\difficulte{3}
\theme{Propriétés de nombres réels, partie entière}
\exo7qcmid{lille-1.reels.018}

\contenu{
\begin{question}
\qclassification{120.03}
\qtags{motcle=partie entière}
On définit la partie fractionnaire d'un réel $x$, par $F(x) = x -E(x)$.
\begin{answers}
    \bad{$F(x) = 0 \iff 0 \le x <1$}

    \bad{Si $7 \le x <8$ alors $F(x) = 7$.}

    \bad{Si $x=-0,2$ alors $F(x) = -0,2$.}

    \good{Si $F(x)=F(y)$ alors $x-y \in \Zz$.}
\end{answers}
\begin{explanations}
La partie fractionnaire est égale à la partie "après la virgule".
Par exemple $F(12,3456) = 0,3456$.
\end{explanations}
\end{question}
}
